% !TeX root = ../main.tex

\chapter{Evaluatie en persoonlijke reflectie}
\label{chap:reflectie}

\section{Evaluatie van het stageplan}
Het oorspronkelijke plan legde de nadruk op dataverzameling, lokale opslag, visualisatie en
voorspelling. De technische invulling die ik in het stageplan had opgeschreven veranderde echter sterk. SQLite werd vervangen door
transparante JSONL-, JSON- en CSV-bestanden. Replay werd uit de eindapp verwijderd en vervangen door
een testgerichte simulator. Een complex AI/ML-traject maakte plaats voor een simpele en uitlegbare
sectorvoorspelling. Tegelijk groeide de race-engineeringfunctionaliteit met meerdere sessiemodi,
pitstopplanning, verschilllende providerondersteuning, meerdere dashboards en grafieken.

Die verschuiving was nodig want binnen zes weken was het belangrijker om de volledige live keten
betrouwbaar te maken dan om een complex model bovenop onzekere data te bouwen. De moeilijkste
problemen zaten niet in de rekenkundige formules, maar in de betekenis en kwaliteit van
timingvelden. Zonder correcte providerinterpretatie zou een geavanceerd model alleen preciezere
foutieve resultaten produceren.

\section{Zelfstandig werken}
De technische uitvoering gebeurde grotendeels zelfstandig.
Zelfstandigheid betekende ook beslissen wanneer een vraag technisch kon worden opgelost en wanneer
domeinfeedback nodig was. De betekenis van een referentietijdregel of geldige pitstop kan niet uit
code worden afgeleid. In zulke gevallen was overleg met de opdrachtgever noodzakelijk voordat
verdere implementatie kon gebeuren.

\section{Communicatie en requirementsanalyse}
Het meermaals overleggen met de klant was een centraal deel van de stage. Presentaties met uitleg en screenshots van wat er was veranderd en
concrete scenario's bleken effectiever dan abstracte technische uitleg. Een vraag als ''Waar staan
we na een pitstop van 75 seconden?'' maakte onmiddellijk zichtbaar welke data en aannames ontbraken.
De implementatie kon vervolgens worden uitgelegd in termen van opeenvolgende verschillen,
rondeachterstanden en klasseposities.

Ik leerde requirements niet als vaste lijst te behandelen. De opdrachtgever kon pas na gebruik
beoordelen of een veld groot genoeg was, of een kleur voldoende opviel en of een vergelijking
tijdens een race werkelijk nuttig was. Tegelijk moest ik ervoor zorgen dat elke wijziging een duidelijke
betekenis hield en niet alleen visueel mogelijk werd.

\section{Software-engineeringvaardigheden}
De stage bracht verschillende onderdelen uit de opleiding samen: modulariteit, datastructuren,
foutafhandeling, testontwerp en user-interfaceontwerp. Vooral \textit{separation of concerns}
werd praktisch belangrijk. Pure shared modules bleken niet alleen beter testbaar, maar maakten ook
meerdere dashboards en grafiekvensters mogelijk zonder logica te kopiëren.
Het project maakte ook het verschil duidelijk tussen syntax en semantiek. Een parser kan een tekst
perfect lezen en toch een fout resultaat produceren wanneer bijvoorbeeld een gap van \code{5L} als vijf seconden wordt
behandeld. Betrouwbare software vereist daarom domeinconstraints en tests met betekenisvolle
voorbeelden.

\section{Onverwachte competenties}

Naast de competenties die ik vooraf had bedacht te verbeteren, ontwikkelde ik tijdens de stage ook vaardigheden waarvan ik niet van verwachtte dat ze zo belangrijk zouden worden.

Een eerste onverwachte competentie was het kritisch beoordelen van externe data. Aan het begin ging ik ervan uit dat waarden op een timingpagina rechtstreeks gebruikt konden worden wanneer ze correct werden ingelezen. Tijdens het ontwikkelen bleek dat een waarde syntactisch geldig kan zijn, maar toch een andere betekenis kan hebben dan verwacht. Ik leerde daarom niet alleen controleren of data aanwezig is, maar ook of de eenheid, context en betekenis ervan correct zijn.

Als tweede leerde ik meer denken in operationele risico's. Software die tijdens een race wordt gebruikt, moet niet alleen correct werken in normale omstandigheden. Ook een wegvallende timingpagina, een foutieve sluiting, ontbrekende sectorinformatie of een veranderende vlagstatus moet voorspelbaar en correct worden afgehandeld. Hierdoor begon ik betrouwbaarheid en herstelbaarheid belangrijker te vinden dan het zo snel mogelijk nieuwe functies proberen toe te voegen.

Als laatste kreeg ik meer inzicht in de volledige oplevering van een softwareproduct. De applicatie moest ook op verschillende besturingssystemen kunnen worden verpakt, geïnstalleerd en bijgewerkt. Daarnaast moesten instellingen en opgeslagen racedata begrijpelijk blijven voor iemand die de interne implementatie niet kent. Deze ervaring leerde mij om niet alleen als ontwikkelaar naar afzonderlijke functies te kijken, maar ook verantwoordelijkheid te nemen voor het uiteindelijke gebruik van het volledige product.

\section{Wat beter kon}
Hoewel het project uiteindelijk een bruikbare applicatie opleverde, zijn er verschillende punten die
achteraf beter of vroeger aangepakt hadden kunnen worden.

\begin{enumerate}
  \item Een formele provider-specificatie aan het begin had sommige interpretatiefouten kunnen
  voorkomen. Vooral het verschil tussen \code{GAP}, \code{DIFF} en \code{INT} bleek afhankelijk van de
  timingprovider en zelfs van de sessie. Omdat de feeds vooral via demo's en screenshots beschikbaar
  waren, werden sommige aannames pas tijdens echte tests ontdekt.

  \item Het concept van ``geldige tempo-data'' had vanaf het begin als aparte domeinregel moeten
  bestaan. Safety-car-ronden, FCY-ronden, rode vlag, pit-inlaps, outlaps en extreme outliers kwamen
  eerst op verschillende plaatsen in de code terug. Daardoor ontstond het risico dat een ronde in de
  ene berekening wel werd uitgesloten en in een andere niet.

  \item De testdata had vroeger dichter bij echte racesituaties mogen liggen. Veel unittests werkten
  met kleine, nette voorbeelden, terwijl echte timingdata ontbrekende sectoren, veranderende
  rijnamen, pitstatussen en rondeachterstanden bevat. De Spa-data werd pas later een belangrijke
  regressiebron. Later werd duidelijk dat racehistoriek op de timingwebsite beschikbaar bleef. Als
  dat vroeger bekend was geweest, had ik sneller met echte racedata kunnen testen.

  \item De ondersteuning voor meerdere gevolgde auto's had vroeger expliciet in het ontwerp kunnen
  zitten. De basisdata hoeft maar één keer opgehaald te worden, maar dashboards, referentietijden,
  instellingen en rapporten zijn wel per auto verschillend. Dat onderscheid werd pas later volledig
  duidelijk.

  \item De configuratie van racespecifieke instellingen had vanaf het begin duidelijker gescheiden
  kunnen worden van live aanpasbare instellingen. Totale raceduur, verplichte pitstops en
  circuitinformatie veranderen normaal niet tijdens een race, terwijl pitstoptijd, referentietijden
  en track condition wel kunnen veranderen. Die scheiding is uiteindelijk toegevoegd, maar had eerder
  voor minder verwarring gezorgd.
\end{enumerate}

